\section{Holographic Indexing: Privacy-Preserving Discovery in Federated Knowledge Networks}

\textbf{Author}: Bryan Alexander Buchorn  
\textbf{Affiliation}: Independent Researcher, Las Vegas, NV, USA  
\textbf{Status}: v2.51.0 (Phase 167 (STRB) COMPLETE)
\textbf{Date}: May 2, 2026

---

\subsubsection{Abstract}
Federated Knowledge Graph reasoning requires nodes to identify relevant information across decen\-tralized peers without compr\-omising data privacy or bandwidth. We present \textbf{Holographic Indexing}, a two-tier discovery protocol designed for "Blind Semantic Search." Our method combines \textbf{Bloom Filter} \cite{bloom1970} membership probing for exact entity matching with \textbf{Community Centroid Signatures} for semantic neigh\-borhood appro\-ximation. This "Holographic" repre\-sentation allows nodes to identify potential reasoning paths in remote peers with minimal data leakage. We formalize the const\-ruction of these signatures and the \textbf{Synaptic Bridge Attention} mechanism that enables cross-graph traversals. We further describe security enhan\-cements in v2.51.0, including \textbf{HMAC-SHA256 Path Provenance} \cite{gentry2009} and \textbf{Federated Leases}, which ensure the integrity and reliability of federated reasoning in high-velocity, multi-tenant envir\-onments. As of v2.51.0, full federated beam execution is realized through \texttt{DistributedBeamTraversal} and a dedicated \texttt{/traverse} endpoint enabling cross-node branch delegation, while the \texttt{ExternalValidator} (Phase 52) extends the federation concept to open scientific literature — validating graph hypotheses against PubMed, arXiv, and OpenAlex.

\subsubsection{1. Introd\-uction}
The expansion of decen\-tralized Knowledge Graphs neces\-sitates a protocol for inter-graph discovery that respects the data sovereignty of individual nodes. Traditional federated search methods often require a central index or the exchange of full node lists, both of which are unacc\-eptable in privacy-sensitive domains (e.g., healthcare or defense). Holographic Indexing addresses this by exchanging compressed, non-reversible topological signatures.

\subsubsection{2. Tier 1: The Entity Hologram}
Each node $\mathcal{G}_i$ generates a Bloom Filter $B_i$ of its entity set $\mathcal{E}_i$. This allows a peer $\mathcal{G}_j$ to verify the existence of a specific entity $e$ with a confi\-gurable false-positive rate $p$, while ensuring that the full set $\mathcal{E}_i$ cannot be enumerated via the filter.

\subsubsection{3. Tier 2: The Community Hologram}
To support fuzzy, semantic discovery, we introduce the Community Centroid Signature. For every community $C_k \in \mathcal{G}_i$, the node computes:
\begin{itemize}
\item \textbf{Centroid}: $\vec{c}_k = \text{mean}(\{\vec{e} : e \in C_k\})$
\item \textbf{Radius}: $\rho_k = \max \|\vec{e} - \vec{c}_k\|$

\end{itemize}
The set of all centroids forms the "Semantic Hologram." A remote reasoning beam looking for concept $\vec{x}$ can compute the "Synaptic Bridge Score":
\begin{equation}
\text{score}(C_k) = \cos(\vec{x}, \vec{c}_k)
\end{equation}
Peers exceeding a threshold $\sigma$ (default 0.75) are flagged as relevant reasoning desti\-nations.

\subsubsection{4. Enterprise Security (v2.51.0)}
To prevent adversarial path injection in federated networks, v2.51.0 implements \textbf{HMAC-SHA256 Path Provenance}. Every reasoning response from a remote adapter is crypt\-ographically signed using a shared secret $\mathcal{K}$. This ensures that "Synaptic Bridge" paths are both seman\-tically valid and struc\-turally authentic.

\subsubsection{5. Conclusion}
Holographic Indexing provides a mathe\-matically robust and privacy-preserving framework for federated graph reasoning. By decoupling exact membership from semantic relevance, it enables secure, decen\-tralized intel\-ligence at scale. In CEREBRUM v2.51.0, full federated beam execution is operational via \texttt{DistributedBeamTraversal} and the \texttt{/traverse} delegation endpoint, while the \texttt{ExternalValidator} extends the federation concept beyond peer KG nodes to open scientific literature, demon\-strating that the Holographic discovery protocol generalizes naturally to heter\-ogeneous federated knowledge sources.

---

\subsection{6. Recent Advances (v2.24.0 -\textgreater  v2.51.0)}

Federated reasoning has been one of the most actively developed areas of the CEREBRUM framework since v2.24.0. The following describes advances directly relevant to this paper.

\textbf{Distri\-butedBeamTraversal and /traverse Endpoint (Phase 32).} The conceptual federation described at v2.24.0 is now fully implemented. \texttt{DistributedBeamTraversal} orche\-strates a multi-node beam search where individual branches can be delegated to remote CEREBRUM nodes via the \texttt{/traverse} REST endpoint. Each remote node runs its own local beam segment and returns ranked partial paths; the coordinator merges and re-ranks these using the global CSA scoring function. This eliminates the need for a central index while preserving the beam's global coherence.

\textbf{RemoteAdapter Archit\-ecture.} A \texttt{RemoteAdapter} implements the standard \texttt{GraphAdapter} interface but routes all graph queries to a remote \texttt{/traverse} endpoint rather than a local data structure. This allows \texttt{DistributedBeamTraversal} to treat remote nodes as transparent graph backends, with Holographic Indexing determining which remote nodes are probed for each query.

\textbf{FederatedAdapter with Procrustes Embedding Alignment.} When multiple remote nodes use indep\-endently trained embedding spaces, cross-node semantic similarity scores are not directly comparable. The \texttt{FederatedAdapter} applies the same Orthogonal Procrustes alignment (PAPER\_008) used for cross-modal signal encoding to align remote embedding spaces to a shared root space before computing Synaptic Bridge Scores. This prevents high-cosine false positives caused by unaligned embedding geometries.

\textbf{ExternalValidator: Scientific Literature Federation (Phase 52).} Beyond peer KG nodes, \texttt{ExternalValidator} extends the federation concept to open scientific databases. Given a hypothesis edge proposed by the \texttt{HypothesisEngine} (PAPER\_007), \texttt{ExternalValidator} queries PubMed, arXiv, and OpenAlex for corro\-borating literature. Validated hypotheses receive elevated confidence scores; contr\-adicted ones are suppressed. This represents a qualitative extension of "federated knowledge discovery" from structured graphs to unstr\-uctured scientific text.

\textbf{Security Enhanc\-ements.} HMAC-SHA256 path provenance, introduced at v2.24.0, is now enforced on all \texttt{/traverse} responses. Federated leases include TTL-bounded tokens to prevent replay attacks in long-running distributed sessions.

---
\subsection{Acknow\-ledgments}

The author gratefully ackno\-wledges the use of Claude (Anthropic) as a research assistant throughout this work. Claude assisted with mathe\-matical forma\-lization, code generation, manuscript preparation, and technical writing. All conceptual contr\-ibutions, archi\-tectural decisions, exper\-imental design, and intel\-lectual claims are solely the author's.

\textbf{References}
\begin{enumerate}
\item Bloom, B. H. (1970). Space/time trade-offs in hash coding with allowable errors. Commun\-ications of the ACM.
\item Gentry, C. (2009). Fully homomorphic encryption using ideal lattices. STOC.
\item Kairouz, P., et al. (2021). Advances and open problems in federated learning. Foundations and Trends in Machine Learning.
\item Broder, A., \& Mitzen\-macher, M. (2004). Network appli\-cations of Bloom filters: A survey. Internet Mathematics.
\item Tarkoma, S., Rothenberg, C. E., \& Lagerspetz, E. (2011). Theory and practice of Bloom filters for distributed systems. IEEE Commun\-ications Surveys \& Tutorials.
\item Buchorn, B. A. (2026). CEREBRUM v2.51: Complete Technical Specif\-ication for Autonomous Knowledge Graph Reasoning. [CEREBRUM\_REPORT\_PLACEHOLDER].

\end{enumerate}
---
\textbf{Reviewed on}: May 2, 2026 for version v2.51.0
